\documentclass[UTF8,fontset=fandol,11pt,a4paper]{ctexart}
\usepackage[a4paper,margin=1.75cm,headheight=15pt]{geometry}
\usepackage{amsmath,amssymb,mathtools,booktabs,tabularx,array,enumitem}
\usepackage[most]{tcolorbox}
\usepackage{tikz}
\usetikzlibrary{arrows.meta,positioning}
\usepackage{xcolor,fancyhdr,hyperref,microtype,lastpage}
\newcolumntype{Y}{>{\raggedright\arraybackslash}X}
\newcolumntype{L}[1]{>{\raggedright\arraybackslash}p{#1}}
\newcolumntype{C}[1]{>{\centering\arraybackslash}p{#1}}
\setlength{\parindent}{2em}
\setlength{\parskip}{0.34em}
\setlength{\emergencystretch}{3em}
\renewcommand{\arraystretch}{1.2}
\setlength{\tabcolsep}{3pt}
\setlist[itemize]{itemsep=0.18em,topsep=0.35em,leftmargin=2em}
\setlist[enumerate]{itemsep=0.18em,topsep=0.35em,leftmargin=2em}
\definecolor{deep}{RGB}{38,58,72}
\definecolor{soft}{RGB}{244,247,248}
\definecolor{greenish}{RGB}{52,91,74}
\definecolor{warnc}{RGB}{136,61,58}
\definecolor{purple}{RGB}{84,72,112}
\definecolor{rulegray}{RGB}{110,116,120}
\hypersetup{colorlinks=true,linkcolor=deep,urlcolor=purple,pdftitle={CO-08 总线与 I/O 硬件}}
\pagestyle{fancy}
\fancyhf{}
\lhead{\small\color{deep}\textbf{408 · 计算机组成原理 CO-08}}
\rhead{\small\color{deep}Request · Transaction · Completion}
\cfoot{\small 第 \thepage\ 页 / 共 \pageref{LastPage} 页}
\renewcommand{\headrulewidth}{0.4pt}
\newtcolorbox{corebox}[1]{breakable,colback=soft,colframe=deep,title=\textbf{#1},boxrule=0.8pt,arc=2mm}
\newtcolorbox{methodbox}[1]{breakable,colback=white,colframe=greenish,title=\textbf{#1},boxrule=0.7pt,arc=1.5mm}
\newtcolorbox{warnbox}[1]{breakable,colback=white,colframe=warnc,title=\textbf{#1},boxrule=0.7pt,arc=1.5mm}
\newcommand{\kw}[1]{\textbf{\color{deep}#1}}
\title{\vspace{-1.1cm}\textbf{总线与 I/O 硬件}\\[0.4em]\large 同步处理器怎样与异步设备完成可确认的协作}
\author{408 · 计算机组成原理 Topic CO-08}
\date{}

\begin{document}
\maketitle

\begin{corebox}{母问题：一次设备工作怎样从意图变成可确认的完成？}
总线、控制器、轮询、中断和 DMA 不是五个孤立名词，而是一条责任链：
\[
\boxed{\text{Request}\to\text{Interface State}\to\text{Arbitration}\to\text{Transaction}\to\text{Transfer}\to\text{Completion Signal}}
\]
每一步都要回答：谁拥有控制权，谁驱动共享资源，数据经过哪里，什么状态证明本步结束。
\end{corebox}

\begin{methodbox}{本册定位与 Owner 边界}
\begin{itemize}
\item \kw{Owns：}总线角色和事务、同步/异步握手、仲裁、有效带宽、控制器寄存器、I/O 编址、polling、中断硬件入口、DMA 控制器和传输模式。
\item \kw{Uses：}CO-02 的 ISA 入口契约、CO-05/06 的内存与 Cache 接口、CO-04 的精确提交边界。
\item \kw{Stops：}驱动策略、请求队列、block/wakeup、buffer policy、DMA mapping/unmapping 和设备调度归 OS/X-B03。
\end{itemize}
\end{methodbox}

\tableofcontents
\newpage

\section{总线是协议，不是一捆静态导线}
共享互连至少包含 initiator/master、target/slave、arbiter 和 interconnect。一次事务通常经过请求、仲裁、地址/命令、一个或多个 data beat、response 与释放。DMA 控制器获授权后也可成为 initiator；“CPU 永远是主设备”不是不变量。

地址、数据与控制信息承担不同责任：地址标识目标，数据线宽度限制每 beat 原始数据量，控制/响应表达 read/write、byte enable、ready、error、request/grant 等协议状态。地址线分时复用不会自动增加数据带宽。

\section{同步、异步与分离事务}
同步总线按共同 clock 采样，慢目标可插入 wait state；异步总线靠 request/acknowledge 握手建立因果完成关系。二者描述\kw{定时协议}，串行/并行描述\kw{物理传输}，不能互推快慢。

分离事务把 request 与 response 解耦，长延迟时可释放通道服务其他请求，但需要事务 ID、缓冲和匹配。是否允许多个 outstanding request 是协议属性，不由“异步”一词推出。

\section{仲裁：选择共享资源的临时 Owner}
仲裁只决定哪个请求者在当前时段获得发起权，不决定 OS 层的设备归属。固定优先级响应关键设备快，却可能饥饿；轮转改善公平性，却增加状态与最坏等待；链式、独立请求或分布式实现还要比较线路成本、仲裁延迟和故障范围。

\begin{center}\small
\begin{tabularx}{\linewidth}{L{2.2cm}YYY}\toprule
机制 & 关键状态 & 优势 & 风险\\\midrule
固定优先级 & pending 与 priority & 关键请求低延迟 & 低优先级饥饿\\
轮转 & last grant / pointer & 长期公平 & 状态与最坏等待增加\\
链式查询 & grant 传播位置 & 线路少 & 延迟累积、物理顺序固化\\
独立请求 & 每主设备 request/grant & 并行判优快 & 线路和逻辑成本高\\
\bottomrule
\end{tabularx}
\end{center}

\section{带宽题先写完整事务}
理论峰值为
\[
BW_{peak}=\frac{bytes}{beat}\times\frac{beats}{cycle}\times frequency.
\]
有效带宽还要扣除仲裁、地址/响应阶段、wait、方向切换、编码与空闲。若一个事务搬运 $D$ 字节、总耗时 $T$，稳定口径是 $BW_{effective}=D/T$。频率、总线周期、数据周期与设备时间必须分开列。

\begin{warnbox}{最小反例}
64 位、1 GHz 只给出每 beat 上限与时钟条件；若每次数据前有地址阶段、目标插入等待，或并非每周期都有 beat，有效带宽就不是自动的 8 GB/s。
\end{warnbox}

\section{I/O 接口把设备变成寄存器状态机}
控制器通常暴露 data/FIFO、status、control/command，以及可选地址、计数、描述符和 interrupt mask。软件写 command 只是启动；busy、ready、done、error 才描述后续状态。

memory-mapped I/O 用普通地址空间和 load/store；isolated I/O 用独立端口空间和专用指令。差别在软件可见寻址契约，不等于设备性能。MMIO 是否可缓存、能否重排和访问宽度须服从 ISA/平台规则。

\section{三种控制方式：谁等待、谁搬数据、谁通知}
\begin{center}\small
\begin{tabularx}{\linewidth}{L{2.0cm}YYY}\toprule
方式 & CPU 等待 & 数据路径 & 完成发现\\\midrule
Polling / PIO & 循环查询状态 & device $\leftrightarrow$ CPU register $\leftrightarrow$ memory & 主动读取 done\\
Interrupt-driven & 可执行其他任务 & 通常仍由 CPU 指令搬运 & 设备异步请求入口\\
DMA & 只配置与收尾 & device $\leftrightarrow$ memory & 块完成/错误通知\\
\bottomrule
\end{tabularx}
\end{center}
中断消除的是持续 busy waiting，不自动消除 ISR 和搬运成本；DMA 下沉的是传输循环，不代表 CPU、Cache、内存和互连没有成本。

\section{中断硬件路径与软件交接}
\[
\text{device event}\to\text{pending}\to\text{mask/priority}\to\text{precise boundary}\to\text{architectural entry state}\to\text{vector}.
\]
外部中断与当前指令异步；异常通常与当前指令同步相关。CPU 是否自动关中断、保存哪些最小状态、向量怎样形成均由 ISA 规定。完整通用寄存器保存、驱动服务、确认设备、唤醒进程和调度属于软件。

中断控制器的 pending、enable/mask、priority 和 in-service 状态必须分开：请求到达不等于立即接受，接受不等于 handler 已完成，handler 返回也不必自动清除设备源头。

\begin{methodbox}{硬件 Stop Boundary}
本 Topic 的最后一个稳定动作是“按 ISA 规定形成可进入软件的体系结构状态”。“唤醒哪个任务、是否重试请求、怎样处理错误”不写成硬件动作。
\end{methodbox}

\section{DMA：把搬运循环下沉到控制器}
CPU 先配置 source、destination、length、direction/control 或 descriptor；DMAC 获得设备和总线条件后发起 transaction，逐 beat 更新地址与计数，完成或错误时置状态并通知 CPU：
\[
\text{Setup}\to\text{DREQ}\to\text{Arbitrate}\to\text{Transfer}\to\text{Count/Address Update}\to\text{Complete}.
\]

burst 一次授权连续传输，吞吐高但占用长；cycle stealing 每次占少量周期，在 CPU 干扰与吞吐间折中；transparent/interleaved 利用约定时隙，依赖具体时序。DREQ 是设备到 DMAC 的传输请求，不是 CPU 中断；总线 grant 也不是上下文切换。

scatter-gather 可由描述符链减少大块非连续传输的重复配置，但描述符建立、地址映射、缓存一致性和完成处理仍产生 CPU/系统成本，因此“DMA 对任意传输恒为 $O(1)$ CPU 开销”不是稳定结论。

\section{贯穿母例：DMA 读取一个块}
\begin{enumerate}
\item 软件准备设备可用的目标地址、长度和方向，并按平台规则处理映射/一致性；
\item CPU 写设备/DMAC 控制寄存器，控制器进入 busy；
\item 设备产生 DREQ，DMAC 作为潜在 initiator 请求互连；
\item arbiter 授权后，DMAC 与内存/设备完成若干 beat，更新地址和剩余计数；
\item 共享内存或总线冲突表现为等待，CPU 是否同时运行取决于其当前资源需求；
\item 计数归零或错误使控制器写 completion state，并产生可屏蔽的 interrupt request；
\item CPU 在精确边界接受中断，形成 ISA 规定的入口状态；
\item 软件确认状态、解除映射、完成请求或唤醒任务。
\end{enumerate}

\section{五列概念边界}
\begin{center}\small
\begin{tabularx}{\linewidth}{L{1.8cm}C{0.35cm}L{1.9cm}YY}\toprule
概念 A & $\neq$ & 概念 B & 真正区别与题面信号 & 混淆后果\\\midrule
request & $\neq$ & transaction & 提出需求 vs 已获权并执行协议 & 漏仲裁与等待\\
bus grant & $\neq$ & interrupt accept & 互连所有权 vs CPU 控制流入口 & 把暂停访存写成 ISR\\
DREQ & $\neq$ & IRQ & 设备请求 DMAC 搬运 vs 控制器通知 CPU & 错画信号方向\\
interrupt entry & $\neq$ & OS service & ISA 最小入口 vs 软件处理与唤醒 & 硬件越界拥有 OS\\
peak bandwidth & $\neq$ & effective bandwidth & 理想 beat 速率 vs 完整事务有效数据率 & 漏协议开销\\
MMIO address & $\neq$ & memory data & 设备寄存器契约 vs 普通可缓存对象 & 错用缓存/重排\\
\bottomrule
\end{tabularx}
\end{center}

\section{做题控制协议}
\begin{enumerate}
\item 标 initiator、target、数据方向、共享资源与 completion 条件；
\item 把 request、arbitration、address/control、data beat、response 分开；
\item 逐阶段写资源 Owner 和允许重叠的 CPU/设备工作；
\item 带宽先算完整事务有效数据/总时间，再与峰值比较；
\item I/O 方式题回答谁等待、谁搬数据、谁通知；
\item 中断题分 pending、accept、entry、software service、return；
\item DMA 题分 setup、transfer、completion，并单列一致性和 OS handoff。
\end{enumerate}

\begin{methodbox}{压缩成第一动作}
看到总线或 I/O 题，先问：\kw{当前谁拥有控制权？当前事件只是 request，还是已经形成 transaction 或 completion evidence？}
\end{methodbox}

\section{全科坐标与跨册交接}
\begin{corebox}{Map Coordinate：Resource / Interface}
CO-08 负责共享线路、控制器状态、仲裁、事务和完成证据。Interrupt 解决“完成后怎样通知 CPU”，DMA 解决“数据怎样批量搬运”；两者可组合但不能互相替代。X-B03/OS-05 接管 mapping、请求队列、completion 和 wakeup。
\end{corebox}
本册 Cost 要把事务有效数据、地址/控制/等待周期、仲裁、方向切换和可重叠阶段写全；总线峰值、有效带宽、设备 latency 与 CPU 指令时间不是同一个量。若访问是 MMIO，属性可能绕过普通 Cache，但仍需回到 CO-07/B02 的权限与地址接口。

\section{考纲映射与来源处理}
\begin{center}\small
\begin{tabularx}{\linewidth}{L{3.0cm}YY}\toprule
考点 & 本册机制 & 接口\\\midrule
总线结构/仲裁/定时 & 角色、事务、握手、有效带宽 & CO-03/05\\
I/O 接口与编址 & 寄存器状态机、MMIO/isolated I/O & CO-02\\
轮询/中断 & 等待、搬运、通知；硬件精确入口 & X-B01/X-B03\\
DMA & setup、仲裁、transfer、completion & OS-05/X-B03\\
\bottomrule
\end{tabularx}
\end{center}

\begin{corebox}{一句话}
总线题追踪“谁在何时拥有共享线路”，I/O 题追踪“谁推进状态、谁搬数据、什么证据证明完成”；中断与 DMA 都只是跨层 handoff，不替 OS 完成软件责任。
\end{corebox}

本册吸收旧《总线结构》《I/O 控制方式》《中断处理》《DMA》及相关控制器卡片；空白《总线计算》《I/O 接口》不产生知识更新，OS 的阻塞、唤醒、buffer policy 与调度保留给各自 Owner。

\end{document}
